Now by
definition we have:
@@M0@@
Furthermore, the exact sequence:
@@M1@@
gives the relation:
@@M2@@
from which we deduce:
@@M3@@

Taking 4.4 into account, the relation to be proved is therefore written:
@@M4@@

Putting @@M5@@, this relation is precisely the Riemann--Roch theorem applied to the regular immersion @@M7@@ of codimension @@M6@@, and to the element @@M8@@ of @@M9@@.

Moreover, since we supposed that @@M10@@ satisfies the hypothesis of (VII 4.1), @@M11@@ is divisible by @@M12@@. One can therefore find @@M13@@ such that @@M14@@, or equivalently @@M15@@ (exp. VII 2.7). We are thus reduced to the case where @@M16@@ is a divisor, and to an element @@M17@@ such that there exists @@M18@@, with @@M19@@.

\sgaStableTarget{sga6.e08.para.4-6}\textbf{4.6.} We must therefore show the relation:
@@M20@@
Now the projection formula gives:
@@M21@@
and from the exact sequence: @@M22@@, where @@M23@@ is invertible here, we obtain:
@@M24@@
hence
